\documentclass{article}
\usepackage{graphicx}
\usepackage{amsmath}
\usepackage{amssymb}
\title{P-01}
\author{A. E. Mahrouss\\amlal@nekernel.org}
\date{\today}

\begin{document}
\maketitle

\section{Abstract}
Example 1.a lays out the weight of a variable $S$ over a frequency $F$ that may be equal to the division of the value between the highest signal and lowest signal of $S$.

\section{Example 1.a}
Let
$\forall S \in \mathbb{R}; \quad Baud(MIN) < S < Baud(MAX)$:
\begin{align*}
Freq(S) = \frac{Hi(S)}{Lo(S)} \cdot Weight(x)\\
-\frac{Freq(S)}{Weight(S)} = \frac{Hi(S)}{Lo(S)}
\end{align*}
By Hartley's Law, the weight of $S$ over its frequency $F$, may be equal of the division between the highest signal and lowest signal of $S$.

\nocite{*}

\bibliographystyle{acm}
\bibliography{ieee_citation}

\end{document}